\Needspace{9\baselineskip}
\originalchapter{信道容量与资源互换}
\originalsection{Holevo信息的含义}
经典标签$X=i$以$p_i$制备$\rho_i$, 总态$\rho_{XB}=\sum_ip_i\ket i\bra i\otimes\rho_i$. 其互信息为$\chi=S(\sum_ip_i\rho_i)-\sum_ip_iS(\rho_i)$. Bob测量得到经典$Y$, 量子互信息在信道$B\to Y$下不增, 所以$I(X:Y)\le\chi$. 该数据处理定理可由量子相对熵单调性得出; 后者是本卷明确引用的高级熵定理, 不以经典log-sum不等式的标量证明冒充完整非对易证明.

同一比特的两正交纯态等权给$\chi=1$且可达. 不正交纯态等权则平均谱为$(1\pm|\langle\psi_0|\psi_1\rangle|)/2$, $\chi<1$. 单份最优可访问信息还可能严格低于$\chi$, 多份联合解码的HSW定理才说明对无记忆经典--量子信道以$\chi$为可达编码率. 平均态典型投影和条件典型投影给随机编码的打包, 检测器可采用平方根测量; 完整非对易打包误差证明作为来源中的高级定理引用, 本卷不虚称已逐行重证.

\originalsection{经典、量子与纠缠辅助容量}
无记忆量子信道传经典信息的一般容量需$\chi$的多次使用正规化最大值, 单次最大$\chi$不总可加. 传量子信息的可达量以相干信息$I_c(\rho,\mathcal N)=S(B)-S(E)=-S(R|B)$组织, 其中$RBE$是输入净化经信道等距后的纯态. 一般量子容量也需正规化且可能有超可加性. 因而不能仅写一个单字母式就称所有信道容量已解决.

预共享无限纠缠辅助传经典信息时,$C_E=\max_\rho I(R:B)$是单字母容量. 一个完美量子比特信道给2经典比特, 与超密编码一致; 一个完美经典比特信道不能仅靠纠缠传未知量子态, 还需传态的两经典比特和一个Bell对. 擦除信道以概率$\epsilon$把量子比特换成正交擦除旗, 可用纠缠与环境对称性得量子容量$\max(0,1-2\epsilon)$, 经典容量$1-\epsilon$, 纠缠辅助经典容量$2(1-\epsilon)$. 这些值依旗信息可见且无记忆的模型.

\originalsection{母协议、父协议与纠缠蒸馏}
对纯$ABE$共享态, 量子状态合并与解耦可描述资源互换: 发送约$\tfrac12I(A:E)$量子比特并提取约$\tfrac12I(A:B)$ ebit是母协议的代表形式, 精确任务还需参考系统与转移方向. 父协议用$\tfrac12I(A:E)$预共享ebit辅助信道, 可传$\tfrac12I(A:B)$量子比特. 二者的核心是让环境与欲传系统解耦, 不仅保证Bob某个边缘态近似正确.

混态可蒸馏纠缠与形成纠缠的成本一般不同, 有PPT束缚纠缠. LOCC不能从可分态制造纠缠; 对纯态熵的平均单调性不能直接延伸成混态$S(\rho_A)$单调性. 本章给资源变量和操作终点, 完整解耦定理、HSW及一般容量编码定理不在本卷的独立证明范围, 作为已核课程笔记结果引用.
